% ======================================================================
\section{Coupling reduced coordinates to learned dynamics}
\label{sec:dynamics}
% ======================================================================

A reduced representation is rarely an end in itself; in most ML applications it
is the state space of a forecasting or surrogate model. This section surveys
how latent dynamics are learned and, more importantly for the choice of DR,
which properties of the latent space make them learnable.

\subsection{Sequential versus joint training}

In \emph{sequential} training the encoder and decoder are fitted first, by POD,
DLS, or an AE trained on reconstruction alone, and the dynamics model is then
trained on the fixed latent trajectories
\parencite{hasegawa2020bluff,nakamura2021channel,solera2024betavae,morton2026scitech}.
This is modular and cheap, and a single representation can serve several
dynamics models. Reconstruction-optimal coordinates, however, are not
necessarily forecasting-friendly. In \emph{joint} training the dynamics loss is
back-propagated into the encoder, shaping the latent space so that its
evolution is simple: linear \parencite{lusch2018koopman,otto2019linearly},
sparse in a function library \parencite{champion2019sindyae}, or accurately
integrable over several steps \parencite{fresca2021dlrom}. Joint training
couples the design choices and is harder to optimize, but it addresses property
(P3) directly. Intermediate strategies fine-tune a pretrained autoencoder with
a small dynamics loss or add temporal-smoothness regularization to the latent
trajectories.

\subsection{Structured regression: operator inference and sparse identification}

When the latent space is linear (POD), physics suggests the structure of the
reduced model. For the incompressible Navier--Stokes equations Galerkin
projection yields a quadratic system
$\dot{\z}=\bm{c}+\bm{A}\z+\bm{H}(\z\otimes\z)$, and operator inference learns
$\bm{c},\bm{A},\bm{H}$ by least squares from projected trajectories
\parencite{peherstorfer2016opinf,kramer2024opinf}. Lifting transformations
extend this to non-polynomial systems \parencite{qian2020liftlearn}, and
quadratic manifolds extend it to curved trial manifolds
\parencite{geelen2023quadratic}. Sparse identification of nonlinear dynamics
(SINDy) \parencite{brunton2016sindy} instead selects a few terms from a library
of candidate functions by sparse regression, yielding interpretable latent
equations. \textcite{loiseau2018sparse} applied it to sensor-derived features.
\textcite{champion2019sindyae} trained an autoencoder jointly with a SINDy model
so that the discovered coordinates admit a parsimonious governing equation.
The LaSDI family extends this to parametric problems: latent ODE coefficients
are identified for each training parameter and interpolated to new ones
\parencite{fries2022lasdi}, with greedy, physics-informed sampling of training
parameters \parencite{he2023glasdi} and Gaussian-process interpolation with
uncertainty quantification \parencite{bonneville2024gplasdi}.

\subsection{Koopman embeddings}
\label{sec:koopman}

Koopman operator theory seeks observables in which a nonlinear system evolves
linearly \parencite{mezic2013,brunton2022koopman}. Deep Koopman methods learn
these observables with an autoencoder and impose linear latent dynamics
$\z_{k+1}=\bm{K}\z_k$ through a multi-step prediction loss
\parencite{takeishi2017koopman,lusch2018koopman,otto2019linearly}.
\textcite{lusch2018koopman} used an auxiliary network to parametrize eigenvalues
as functions of the state, capturing systems with continuous spectra.
\textcite{morton2018deep} learned such models of cylinder-wake dynamics and
used them for model-predictive control, and \textcite{pan2020physics}
introduced a probabilistic formulation with guaranteed stability. For a
low-order model of near-wall turbulence, \textcite{eivazi2021koopman} found that
a Koopman framework with nonlinear forcing reproduced the long-term statistics
as accurately as an LSTM at lower cost, and gave better short-term
predictions. Koopman embeddings also serve as tokenizers for transformers
\parencite{geneva2022transformers}.

\subsection{Recurrent networks, reservoir computing, and neural ODEs}

Recurrent networks, especially LSTMs \parencite{hochreiter1997lstm}, were the
first widely used black-box latent dynamics models.
\textcite{vlachas2018lstm} forecast high-dimensional chaotic systems with LSTMs
in a reduced-order space. For the nine-equation model of turbulent shear flow,
\textcite{srinivasan2019} reproduced turbulence statistics with sub-percent
errors using LSTMs. The LSTM-on-POD and LSTM-on-CAE combinations of
\cref{sec:pod,sec:autoencoders} follow the same pattern
\parencite{wang2018model,mohan2018lstm,eivazi2020deep,hasegawa2020bluff,maulik2021advection,borrelli2022channels}.
Reservoir computing trains only a linear readout of a large random recurrent
network. It forecasts spatiotemporal chaos well when combined with spatial
localization \parencite{pathak2018reservoir} or with a CAE latent space
\parencite{racca2023caesn}, and \textcite{vlachas2020backprop} compared it
systematically with back-propagated RNNs. Neural ODEs \parencite{chen2018node}
learn a continuous-time vector field $\dot{\z}=\bm{f}_{\theta}(\z)$, which
suits irregularly sampled data and connects naturally with dynamical-systems
analysis. In autoencoder coordinates of the Kuramoto--Sivashinsky equation,
\textcite{linot2022node} found that dimension reduction improves performance
relative to learning in the ambient space, and identified a ``sweet spot''
in latent dimension. Stabilized variants add a learned linear damping term for
long-time forecasting \parencite{linot2023stabilized}, and the approach has been
applied to minimal-flow-unit Couette turbulence \parencite{linot2023couette}.
Learning the effective dynamics (LED) \parencite{vlachas2022led} couples an AE
and a recurrent latent propagator with intermittent fine-scale simulation in
an equation-free framework, reducing computational effort by up to two orders
of magnitude while maintaining accuracy.

\subsection{Attention-based sequence models}

Transformers \parencite{vaswani2017attention} replace recurrence with
self-attention over a context window of past latent states, which captures
long-range temporal dependencies and trains in parallel.
\textcite{wu2022transformer} combined a transformer with POD for flows past a
cylinder and a group of buildings, and \textcite{geneva2022transformers} with
Koopman embeddings for several physical systems. \textcite{hemmasian2023transformers}
used a CAE latent space and reported competitive results even for turbulent
datasets. \textcite{solera2024betavae} and \textcite{wang2024optimal} used
$\beta$-VAE latent spaces. Transformers have also been used without a
separate autoencoder, on subsampled or pre-processed fields:
\textcite{yousif2023transformer} generated turbulent inflow planes for
spatially developing boundary layers by combining a transformer with a
super-resolution GAN, and \textcite{sarma2024tbl} predicted actuated turbulent
boundary-layer fields after a dimensionality-reducing pre-processing step.
\textcite{morton2026scitech} paired transformers with DLS latent spaces for
cavity flows at $\mathrm{Re}=15{,}000$ and $30{,}000$
(\cref{sec:dls}). Attention can also operate over \emph{space}: when the latent
state is a set of spatially anchored tokens, as in DLS, AROMA, or universal
physics transformers
\parencite{deshmukh2020dls,serrano2024aroma,alkin2024upt}, the sequence model can
attend across tokens as well as time, and forecasting itself can be posed
locally.

\subsection{Latent generative models}

Deterministic models trained with a mean-squared loss converge to the
conditional mean of future states. For chaotic flows this is smoother and less
energetic than any single realization, and the resulting under-dispersion
often drives rollout drift. Generative models address this by learning the
distribution of future states. Latent diffusion models, developed for images
\parencite{rombach2022ldm}, have been adapted to turbulence generation with
conditional-neural-field latents \parencite{du2024cnf} and to Bayesian
conditional generation \parencite{gao2024bayesian}. \textcite{kohl2023acdm}
benchmarked autoregressive conditional diffusion models for turbulent flow
prediction. Even simple diffusion approaches matched or outperformed
established methods in accuracy and temporal stability, at a higher inference
cost, while providing samples consistent with the flow statistics. PDE-Refiner
\parencite{lippe2023pderefiner} achieves accurate long rollouts through an
iterative refinement that attends to low-amplitude frequency components. The
compression provided by DR is especially valuable for diffusion models, whose
cost scales with the dimension of the space in which they denoise.

\subsection{What makes a latent space good for dynamics?}
\label{sec:goodlatent}

Across these studies, several latent-space properties recur as predictors of
successful forecasting, beyond reconstruction accuracy.
\begin{itemize}[leftmargin=1.5em,itemsep=2pt]
  \item \emph{Sufficiency and Markovianity.} The latent state must retain
  enough information to determine its own future. Aggressive truncation removes
  small scales that feed back on large ones, so the latent dynamics become
  non-Markovian. This is the origin of the closure problem in POD--Galerkin
  models \parencite{wang2012closure,ahmed2021closures}, and it appears in
  learned models as the need for long input histories or memory terms
  \parencite{maulik2020timeseries}.
  \item \emph{Smoothness and low temporal complexity.} Latent coordinates that
  vary smoothly and with few dominant frequencies are easier to forecast.
  Near-orthogonal, disentangled coordinates
  \parencite{eivazi2022betavae,solera2024betavae} and ordered coordinates
  \parencite{fukami2020hierarchical} were motivated partly by this. Localized
  coordinates may reduce temporal complexity in a different way: each local
  coefficient reflects only nearby structures, so its time series can carry a
  narrower range of scales than a global mode that mixes contributions from
  the whole domain. This is one motivation for the localized latent spaces of
  \textcite{morton2026scitech}.
  \item \emph{Appropriate dimension.} Too small a latent space discards
  dynamics, and too large a space adds noisy, weakly constrained directions in
  which errors grow; \textcite{linot2022node} found an intermediate optimum.
  \item \emph{Stability.} Latent spaces or dynamics models with built-in
  stability \parencite{pan2020physics,linot2023stabilized}, and training on
  multi-step rollouts instead of single steps
  \parencite{lusch2018koopman,kohl2023acdm}, reduce error accumulation.
  \item \emph{Statistical fidelity.} For turbulence, matching second-order
  statistics and spectra matters as much as pointwise accuracy, which favours
  representations that retain small-scale energy
  \parencite{deshmukh2018energy,glaws2020insitu} and training objectives or
  generative models that do not collapse to the conditional mean
  \parencite{kohl2023acdm,du2024cnf}.
\end{itemize}
